% TOP_PROBLEM: {"id": 3356, "title": "3 is a primitive root modulo primes of the form 16 q^4 + 1, where q>3 is prime", "category_id": 1, "category": {"id": 1, "name": "number_theory", "display_name": "Number Theory", "description": "Properties of integers, prime numbers, Diophantine equations.", "slug": "number-theory", "order_index": 1, "created_at": "2026-07-31T15:26:25.670Z"}, "status": "open", "problem_number": "OPG-37396", "set_id": 12}
% FIRST_POSTED: 2026-10-01
% ENTRY_KIND: statement_only
% UnsolvedMath ID 3356; source catalogue code OPG-37396
% !TeX program = lualatex
% Definition and sources only; this file is not an LLM solution attempt.
\documentclass[11pt,a4paper]{article}
\usepackage[margin=25mm]{geometry}
\usepackage{fontspec}
\setmainfont{lmroman10-regular.otf}[BoldFont=lmroman10-bold.otf,
 ItalicFont=lmroman10-italic.otf,BoldItalicFont=lmroman10-bolditalic.otf]
\usepackage{amsmath,amssymb,mathtools}
\usepackage{unicode-math}
\setmathfont{latinmodern-math.otf}
\AtBeginDocument{\let\setminus\smallsetminus}
\usepackage{xurl,hyperref}
\hypersetup{colorlinks=true,urlcolor=blue}
\setcounter{secnumdepth}{0}
\setlength{\parindent}{0pt}
\setlength{\parskip}{5pt}
\setlength{\emergencystretch}{3em}
\begin{document}
\section*{3 is a primitive root modulo primes of the form $16q^4+1$, where $q>3$ is prime}\label{3356}
{\small UnsolvedMath ID 3356 · OPG-37396 · Number Theory · OpenGarden}
\subsection{Definitions and mathematical statement}
For a prime $p$ and an integer $a$ not divisible by $p$, the
multiplicative order $\operatorname{ord}_p(a)$ is the least positive
integer $r$ such that $a^r\equiv1\pmod p$. A primitive root modulo
$p$ is an element of order $p-1$, so its powers generate the entire
group $(\mathbb Z/p\mathbb Z)^\times$.

The source's conjecture concerns the thin family
\[
 q>3\text{ prime},\qquad p=16q^4+1\text{ prime}:
               \quad \operatorname{ord}_p(3)=16q^4.
\]
Both primality assumptions are hypotheses. No assertion that this
family contains infinitely many primes $p$ is required. Since the only
prime divisors of $p-1$ are $2$ and $q$, the order condition is
equivalent to the simultaneous noncongruences
\[
 3^{8q^4}\not\equiv1\pmod p,
 \qquad 3^{16q^3}\not\equiv1\pmod p.
\]
Indeed a proper divisor of $p-1$ divides $(p-1)/r$ for some prime
divisor $r$ of $p-1$. This equivalent finite criterion fixes exactly
what must be verified for a proposed example. Prove the assertion for
every admissible $q$, or exhibit primes $q,p$ of the specified form
for which the order is smaller.
\subsection{Short English statement}
Whenever q and 16q to the fourth power plus one are both prime, with q greater than three, must three generate all nonzero residues modulo the latter prime?
\subsection{Sources}
\begin{itemize}
\item[S1] ULAM.AI and the UnsolvedMath Contributors, supplied problems.json, record 3356 (OPG-37396), OpenGarden collection. Catalogue identity and source transcription; CC BY 4.0.
\url{https://huggingface.co/datasets/ulamai/UnsolvedMath}.
\item[S2] Open Problem Garden, 3 is a primitive root modulo primes of the form 16 q\textasciicircum{}4 + 1, where q>3 is prime. Original problem and discussion; the mathematical formulation below is expanded from the supplied source record.
\url{http://www.openproblemgarden.org/op/primes_p_such_that_3_is_a_primitive_root_modulo_p}.
\end{itemize}
\end{document}
